\thispagestyle{empty}
\chapter*{译注本说明}
\addcontentsline{toc}{chapter}{译注本说明}

\textbf{底本。}本书是 S. P. Novikov 主编的《拓扑学 I：总论》（\emph{Topology I: General Survey}），
属 Springer《数学百科全书》（Encyclopaedia of Mathematical Sciences）第 12 卷，
英译本由 B. Botvinnik 与 R. Burns 自俄文译出（1996 年）。中文译注本以科学出版社影印本为工作底本，
逐页对照原书扫描图，正文页码一律标注原书印刷页码，读者可随时回查原书。

\textbf{内容。}全书依次为：导言；第 1 章 最简单的拓扑性质；第 2 章 拓扑空间、纤维化与同伦；
第 3 章 单纯复形与 CW 复形、同调与上同调，及其与同伦论的关系和障碍理论；
第 4 章 光滑流形；结束语；附录 A 3 维流形与纽结拓扑学的新进展；
附录 B 术语对照表；附录 C 原书排印问题一览；参考文献；索引。
其中附录 B、附录 C 与本译注本的全部译注、读书笔记同为译注本所增，非原书文字。

\textbf{译注体例。}
\begin{itemize}[leftmargin=2.2em,itemsep=2pt]
\item 每节标题后附「本节重点」一栏，一至三句，点出本节的核心对象、工具与结论。
\item 每章（含附录）末尾附「本章读书笔记」，综述该章的脉络、核心结论与阅读难点。
\item 正文中夹注 \texttt{Tutor Note}（排为脚注），提示符号约定、省略的论证与背景直觉；
      原书排版有误或术语歧义处，正文照原书，注中说明。
\item 原书公式与定理编号一律照录，不作改动；每至原书翻页处，页边以灰色小字标出对应的原书页码（如 p.\,38），
      读者可据此逐页回查原书扫描本。
\end{itemize}

\textbf{术语。}人名保留英文原形。术语以书末《术语对照表》为准；
首次出现的术语以「中文（English）」形式给出。

\textbf{分工。}本译注本由多个译者分章翻译，另设核验者逐章比对原书页图；
凡与页图不符处均经修订。疏漏之处，敬请读者指正。
\clearpage
